% ====================================================================
% si_fr.tex — [SI-FR] Si-host 정칙용액(Frumkin, Ω<2RT 고용체) 커널 (신규 소절)
%   삽입 위치: Chapter 3 Si, §3.2 케이스 소절(ch3v22_sec02_cases, \label{sec:si-cases}) 인접.
%   저작 sub W5 초안 — 물리 기작·직관 전면(WHY a-Si 는 sharp peak 없는 고용체). 표기·라벨은 기존 계승.
%   저작 근거·가정·불확실점·신규 bib 서지는 같은 폴더 NOTES.md.
% ====================================================================
\subsection{Si-host 정칙용액 커널 --- a-Si 는 왜 sharp peak 없는 단일상 고용체인가}\label{ssec:si-frumkin}
케이스 소절(\S\ref{ssec:si-elemental})이 순환 a-Si 의 dQ/dV 가 plateau 없는 \emph{완만한 경사}로 옮겨감을
실측으로 보였다. 이 소절은 그 개형을 낳는 \emph{단일상 고용체} 물리를 앞세워, Si-host 전이에 들어갈 평형 커널을
정규용액(Frumkin) 형으로 닫는다. 흑연 두-상이 평형에서 \emph{델타}이고 broadening 이 종을 만든 것(\S\ref{sec:broadening})과
정반대로, a-Si 는 \emph{평형 자체가 넓은 종}이라 broadening 물음이 생기지 않는다 --- 그 까닭이 커널의 부호 하나에
새겨진다.

\textbf{(a) 왜 a-Si 는 단일상인가 --- 무질서 host 의 자리 에너지 분포.} 결정질 상은 Li 자리가 하나의 잘 정의된
에너지를 갖고 그 자리들이 협동적으로 채워지며 초격자(예: 흑연 $\sqrt3\!\times\!\sqrt3$)로 \emph{정렬}한다 --- 이 협동성이
miscibility gap(두-상)의 근원이다. 비정질 Si 는 그 반대다: 무질서한 Si 골격은 Li 자리의 \emph{연속적 에너지 분포}를
주고, 장거리 정렬할 초격자가 없어 협동적 상전이가 일어나지 않는다. 게다가 삽입된 Li 사이 유효 상호작용이 무질서
매질에서 약하게 가려져(screened) 정규용액 인력이 문턱 아래($\Omega<2RT$)에 머문다. 두 사실이 합쳐져 a-Si 는
\emph{하나의 상 안에서 조성이 연속으로 변하는 고용체}이고, 전위-조성 곡선 $U(x)$ 가 plateau 없이 매끄럽게
기운다 --- 제일원리가 이 경사 등온선을 재현하고\cite{chevrier_dahn2009}, 비정질 $\mathrm{Li}_x\mathrm{Si}$ 의
제일원리 상도표가 이산 이차상 없는 연속 조성 공간을 보인다\cite{artrith2018}.

이 단일상 성격은 \emph{폭}으로 정량 진단된다. 두-상 전이는 평형이 델타라 폭$/(RT/F)\!\ll\!0.12$(거의 0)인 반면,
a-Si 전이의 실측 폭은 오히려 $RT/F$ 보다 넓다 --- 본 세션 정칙용액 적합(regsol\_si)에서 세 Si-host 전이의
폭$/(RT/F)=[1.45,\,2.74,\,1.09]\gtrsim1$ 로, 흑연 두-상 대비 \emph{$\sim$20--50배} 넓다(내부 수치 검증, tier B).
곧 ``폭이 $RT/F$ 규모거나 그 이상''이라는 것 자체가 두-상이 아니라 고용체라는 실측 판정이다. 예외는 \emph{하나}뿐이다
--- 결정질 $c$-$\mathrm{Li_{15}Si_4}$ 는 첫 리튬화 깊은 끝($\sim$50--60 mV)에서 결정화하며 유일한 두-상 plateau 를
남기나\cite{obrovac_christensen2004,limthongkul2003}, 이는 \emph{1차} 리튬화 feature 라 순환 a-Si 에는
부재한다(첫 사이클만 두-상, 이후 고용체 경로\cite{wang_asi2013,mcdowell_coreshell2013}). 곧 개형이 사이클 이력에
의존한다는 \S\ref{ssec:si-elemental} 의 관측이 이 상 성격 전환의 다른 얼굴이다.

\textbf{(b) Frumkin 단일상 peak 커널 --- 유도.} 커널은 흑연과 같은 정규용액 전위에서 나온다. 전이의 점유 $\theta$
(host Li 함량)에 대한 평형 전위는 식~\eqref{eq:mu} 의 화학퍼텐셜을 $V=(\mu^\circ-\mu)/F+\text{const}$ 로 읽어
\begin{equation}
V(\theta)=U^\circ-\frac{RT}{F}\ln\frac{\theta}{1-\theta}-\frac{\Omega}{F}(1-2\theta)
\label{eq:si-fr-V}
\end{equation}
(흑연 식~\eqref{eq:lco-Veq} 를 여집합 좌표 $\xi=1-\theta$ 로 옮긴 것과 동일 대수). 이를 $\theta$ 로 미분하면
\begin{equation}
\frac{\dd V}{\dd\theta}=-\frac1F\!\left[\frac{RT}{\theta(1-\theta)}-2\Omega\right],
\label{eq:si-fr-dVdtheta}
\end{equation}
곧 식~\eqref{eq:gr-dmudtheta} 의 반충전 곡률과 같은 괄호다(부호 반대 --- $V$ 가 $\mu$ 의 음화). 전이 용량 $Q$ 가
$\theta\!:\!0\!\to\!1$ 로 실려($\dd Q/\dd\theta=Q$), 미분용량은 연쇄율 $\dd Q/\dd V=(\dd Q/\dd\theta)/(\dd V/\dd\theta)$ 로
\begin{equation}
\boxed{\;\Big(\frac{\dd Q}{\dd V}\Big)_{\mathrm{Si}}
=\frac{Q\,F}{\big|\,RT/[\theta(1-\theta)]-2\Omega\,\big|}\,,\qquad \Omega<2RT\ (\text{단일상 고용체})\;}
\label{eq:si-fr-kernel}
\end{equation}
--- 이것이 Si-host 전이의 Frumkin 단일상 peak 커널이다\cite{leviaurbach1999}. 분모가 $0$ 이 되지 않는($\Omega<2RT$)
한 유한 종이며, 두-상 커널처럼 델타로 붕괴하지 않는다.

\textbf{(c) 커널의 거동 --- sharpen$\cdot$broaden 과 폭 이중지위.} 반충전 $\theta=\tfrac12$ 에서 분모는
$|4RT-2\Omega|$ 라, 순높이 $(\dd Q/\dd V)_{\max}=QF/|4RT-2\Omega|=[Q/(4w)]/|1-\Omega/2RT|$($w=RT/F$)다. 세
극한이 물리를 준다:
\begin{itemize}[leftmargin=1.5em,itemsep=1pt]
\item $\Omega\!\to\!2RT^-$: 분모 $\to0$ 이라 peak 가 \emph{좁아지고 발산}한다 --- 유효 폭
$w^{\eff}=(RT/F)(1-\Omega/2RT)\!\to\!0$. 이것이 두-상 문턱을 향한 sharpening 이며, $\Omega\!\ge\!2RT$ 를 넘는 순간
평형이 델타(두-상)로 넘어간다(\S\ref{sec:gr-2L} 식~\eqref{eq:gr-criterion}).
\item $\Omega=0$(이상 격자기체): $w^{\eff}=RT/F$ 의 logistic 미분 종(식~\eqref{eq:eqpeak})으로 정확히 환원 ---
Si-host 미지정 시의 로지스틱 폴백이 이 극한이다.
\item $\Omega<0$(유효 반발): 분모가 $4RT+2|\Omega|$ 로 커져 peak 가 \emph{넓어진다} --- 폭$/(RT/F)>1$ 의 실측(위
(a))과 부합하며, 넓은 galleries 를 $\Omega<0$ 또는 다중도 $n>1$ 로 담는다.
\end{itemize}
여기서 앞 소절의 이중지위(\S\ref{sec:width-w})가 그대로 살아난다 --- a-Si 는 $\Omega<2RT$ 의 단상 쪽에 확고히
놓이므로 그 폭 $w^{\eff}=(RT/F)(1-\Omega/2RT)$(또는 $n>1$ 로 넓힌 폭)는 \emph{평형이 예측하는 검증 가능한 폭}
(첫째 지위)이지 두-상의 현상학적 자유 폭이 아니다. 흑연 두-상과 Si 고용체가 \emph{같은} 정규용액 커널의 두 극한이라는
것이 이 소절의 통일점이다.
\begin{verifybox}
검산 --- (i) 폭 진단: 두-상 폭$/(RT/F)\!\ll\!0.12$ 대 a-Si $[1.45,2.74,1.09]$ 는 식~\eqref{eq:si-fr-kernel} 의
$|1-\Omega/2RT|$ 인자로 $\sim$20--50배 차 --- 부호($\Omega\!\gtrless\!2RT$)가 폭 척도를 가른다. (ii) 면적 보존:
$\int(\dd Q/\dd V)_{\mathrm{Si}}\dd V=Q\!\int_0^1\dd\theta=Q$ 로 전이 총 전하 --- $\Omega$ 값과 무관(sharpen/broaden
은 높이$\times$폭 재분배일 뿐 면적 불변). (iii) 폴백 bit-exact: $\Omega\!\to\!0$ 에서 식~\eqref{eq:si-fr-kernel} 이
로지스틱 종~\eqref{eq:eqpeak} 과 글자까지 일치 --- 코드 \code{kernel:'regsol'} 미지정 시 로지스틱 복귀의 이론짝.
(iv) 단일 두-상 예외: $c$-$\mathrm{Li_{15}Si_4}$($\sim$50--60 mV, 1차 리튬화)만 $\Omega\!>\!2RT$ 측, 순환 a-Si 엔
부재라 커널은 단상형만 쓴다.
\end{verifybox}

\textbf{(d) Si-host $=$ 소수 broad gallery, 그리고 블렌드 가산 중첩.} 연속 고용체를 forward 모델로 담는 방식은
그 연속 등온선을 \emph{소수의 넓은 gallery 전이}로 이산화하는 것이다 --- 각 gallery 가 식~\eqref{eq:si-fr-kernel} 의
넓은 종 하나이고, 이들의 합이 매끄러운 $U(x)$ 를 근사한다(다반응$\cdot$speciation 의 전하 보존 구조\cite{verbrugge_lisi2016},
$\Omega$/$\omega$ 정규용액 형식\cite{msmr_origin2017}). 이 Si-host 커널은 혼합 음극 합성식
~\eqref{eq:blend-dqdv} 의 Si 이중합 몫에 그대로 꽂힌다 --- 두 host 가 \emph{같은 전위 손잡이}(공통 $\mu$)를
공유하므로, 관측 dQ/dV 는 흑연 종과 Si Frumkin 종의 \emph{같은 전위축 위 선형 합}이다:
\begin{equation}
\frac{\dd Q}{\dd V}\bigg|_{U_\mathrm{oc}}=C_\bg+\sum_{\mathrm{host}}\sum_j
\frac{Q_j^{\mathrm{host}}\,\xi_{\eq,j}^{\mathrm{host}}(1-\xi_{\eq,j}^{\mathrm{host}})}{w_j^{\mathrm{host}}}
\bigg|_{\mathrm{Si\ 몫\ \to\ 식~\eqref{eq:si-fr-kernel}}},
\label{eq:si-fr-blend}
\end{equation}
곧 Si 몫만 로지스틱 종에서 Frumkin 종으로 바꾼 것이다. 블렌드 실측이 두 host 기여의 중첩으로 읽힌다는 것
(``clearly a superposition'')은 축소차수 블렌드 모형이 직접 보고하며\cite{tu_blend2024}, 저-Si 복합에서 Si$\cdot$흑연
peak 가 부분 분리로 관측되는 것(\S\ref{ssec:si-sic}\cite{naboka_sic2021})이 그 실측 접점이다 --- 단 이 가산 중첩은
\emph{평형(공통 $\mu$)에서 정확}하고 유한 율속 편차는 GS-2 공백(\S\ref{ssec:si-blend-gs2})이 별도로 진단한다.

\begin{keybox}
\textbf{Si-host 커널 한 줄.} a-Si 는 무질서 host 의 자리 에너지 분포 $+$ 약한($\Omega<2RT$) Li--Li 상호작용이라
\emph{단일상 고용체}이고, 폭$/(RT/F)\!\gtrsim\!1$(실측 $[1.45,2.74,1.09]$; 흑연 두-상 $\ll0.12$ 의 20--50배)이 그
증거다. 평형 커널은 Frumkin 단일상 peak
$(\dd Q/\dd V)_{\mathrm{Si}}=QF/|RT/(\theta(1-\theta))-2\Omega|$(식~\eqref{eq:si-fr-kernel}, $\Omega<2RT$)로,
$\Omega\!\to\!2RT$ 서 sharpen$\cdot$발산 $\cdot$$\Omega<0$ 서 broaden $\cdot$$\Omega\!\to\!0$ 서 로지스틱 폴백
(bit-exact)한다 --- 흑연 두-상과 \emph{같은 정규용액 커널의 두 극한}. 유일 두-상은 $c$-$\mathrm{Li_{15}Si_4}$
(1차 리튬화 $\sim$50--60 mV, 순환 a-Si 부재). Si-host 는 소수 broad gallery 로 이산화되어 공통-$\mu$ 합성
~\eqref{eq:si-fr-blend} 에 가산 중첩된다\cite{tu_blend2024}.
\end{keybox}

% 자산(comp_R1 W5 신규): [W5-SIFR-01 a-Si 단일상 물리 기작(무질서 자리 에너지 분포·Ω<2RT screened)]
%   [W5-SIFR-02 폭 진단 [1.45,2.74,1.09]≳1 vs 흑연 ≪0.12(20–50×)] [W5-SIFR-03 Frumkin 커널 유도 boxed eq:si-fr-kernel(eq:mu·eq:si-fr-V·eq:si-fr-dVdtheta)]
%   [W5-SIFR-04 3극한 sharpen/폴백/broaden·폭 이중지위 단상 첫째 지위] [W5-SIFR-05 verifybox 폭진단·면적보존·폴백 bit-exact·Li15Si4 예외]
%   [W5-SIFR-06 소수 broad gallery 이산화·공통-μ 가산중첩 eq:si-fr-blend(eq:blend-dqdv Si몫 교체)·tu2024 superposition]
